\chapter{韵母}

\section{元音音位}

米亚罗话共有6个元音音位，如表\ref{vowel}所示。

\begin{table}[H]
\vspace{5pt}
    \centering
    \caption{元音音位}
    \label{vowel}
\begin{tabular}{ccc}
   \toprule
    元音 & 例子 & 释义 \\
   \midrule
   /a/ & [\textbf{kɐ}-] & “动词词典形前缀”  \\
   /e/ & [kɐ\textbf{le}] & “兔子”  \\
   /i/ & [tɐ-\textbf{mî}] & “脚”  \\
   /ə/ & [kə-\textbf{lə}] & “虫子”  \\
   /o/ & [\textbf{ko-}] & “向上游（趋向前缀）”  \\
   /u/ & [tɐ-\textbf{ku}] & “头”  \\
   \bottomrule
\end{tabular}
\end{table}

各个元音在在辅音/t/后的准最小对如表\ref{vowelcon}。

\begin{table}[H]
\vspace{5pt}
    \centering
    \caption{元音准最小对}
    \label{vowelcon}
\begin{tabular}{ccc}
   \toprule
    元音 & 例子 & 释义 \\
   \midrule
   /a/ & [kɐ-\textbf{tɐ̂}] & “蒸”  \\
   /e/ & [\textbf{te}-] & “一（数量前缀）”  \\
   /i/ & [ŋɐ-\textbf{tî}] & “（我的）大叔/大舅”  \\
   /ə/ & [\textbf{tə̂}] & “向下游（adv.）”  \\
   /o/ & [\textbf{to}] & “上（n.）”  \\
   /u/ & [\textbf{tû}] & “向上（adv.）”  \\
   \bottomrule
\end{tabular}
\end{table}

音位/a/存在音位变体[ɐ]、[a]、[ɑ]，其中[a]出现在韵尾/n/之前、[ɑ]出现在韵尾/ŋ/之前，其他情况下均为[ɐ]。例如[sman]“药”、[mɑŋmə̂]“军队”、[tə-smɐ̂t]“下半身”。

音位/ə/存在音位变体[ə]、[ɤ]、[ɿ]。出现在软腭音的开音节词尾时为[ɤ]，例如[skɤ̂]“佛”、[tekʰɤ̂]“一万”、[tɑ-ŋɡɤ̂]“背部”、[kɐ-rŋɤ̂]“蹲”。出现在齿龈、龈后的擦音、塞擦音（即/s, z, ʃ, ʒ, ts, tsʰ, dz, tʃ, tʃʰ, dʒ/）后的开音节词尾时为[ɿ]，例如[tə-tʃɿ]“水”、[kə-tsɿ]“猴子”、[tsʰɿ]“野葱”、[tɐ-ʃɿ]“血”、[kɐ-ʒɿ]“告状”、[tə-dʒɿ̂]“皮肤”、[kɐ-zɿ]“管教”、[kɐ-sɿ]“倚靠”等。

\section{韵尾}

在韵尾位置可以出现的辅音为/m, n, ŋ, p, t, k, s, r, w, j, ɲ, ʃ/，其中/ɲ, ʃ/仅作为形态音位出现。韵尾位置的辅音清浊中和。词干中能够出现的韵尾如表\ref{rhyme}所示。

\begin{table}[H]
\vspace{5pt}
    \centering
    \caption{韵母表}
    \label{rhyme}
    \resizebox{\textwidth}{!}{
        \begin{tabular}{c|cccccccccccccc}
   \toprule
	 & m & n & ŋ & p & t & k & s & r & w & j & ms & ks & ps & ŋs \\
    \midrule
	a & /am/ & /an/ & /aŋ/ & /ap/ & /at/ & /ak/ & /as/ & /ar/ & /aw/ & /aj/ & /ams/ & /aks/ & /aps/ &  \\
	e & /em/ & /en/ &  & /ep/ & /et/ & /ek/ & /es/ & /er/ &  & /ej/ & /ems/ &  & /eps/ &   \\
    ə & /əm/ & /ən/ & /əŋ/ & /əp/ & /ət/ & /ək/ & /əs/ & /ər/ & /əw/ &  &  &  &  & /əŋs/  \\
	o & /om/ & /on/ & /oŋ/ & /op/ & /ot/ & /ok/ & /os/ & /or/ &  &  & /oms/ & /oks/ &  & /oŋs/  \\
	u &  & /un/ & /uŋ/ & /up/ & /ut/ & /uk/ & /us/ & /ur/ &  & /uj/ &  &  &  &   \\
	i & /im/ & /in/ & /iŋ/ & /ip/ & /it/ & /ik/ & /is/ & /ir/ &  &  &  & /iks/ & /ips/ &   \\
    \bottomrule
\end{tabular}
    }
\end{table}

对于以/j/、/w/为韵尾的韵母，不将韵母视为双元音，而视为韵核加上辅音韵尾的理由是其在音系过程中展现出与其它辅音韵尾一致的行为。正如在\ref{reduplication}节看到的，部分重叠过程中滋生音节会删除原音节的韵尾部分，而/j/、/w/同样会被删除。例如：

\begin{exe}
\ex kəktsej “小的/年轻的” $\rightarrow$ kəktsektsej \\
\gll kə- ktse- ktsej \\
     $\sigma_{1}$ $\sigma_{2}'$ $\sigma_{2}$ \\
\end{exe}

\begin{exe}
\ex kəxɐw “好的” $\rightarrow$ kəxɐxɐw \\
\gll kə- xɐ- xɐw \\
     $\sigma_{1}$ $\sigma_{2}'$ $\sigma_{2}$ \\
\end{exe}



\subsection{在韵尾位置的复辅音}

米亚罗话在词干的韵尾位置允许出现复辅音，同时一些形态学过程也会产生复辅音韵尾。

米亚罗话词干中能够出现的复辅音韵尾有/ms/、/ks/、/ps/、/ŋs/，例如表\ref{comcoda}所示。

\begin{table}[H]
\vspace{5pt}
    \centering
    \caption{韵尾复辅音示例}
    \label{comcoda}
\begin{tabular}{c|cc}
     \toprule
复辅音   & 示例              & 释义\\
     \midrule
	-ks & {[}zgrôks{]}      & 手镯    \\
-ps & {[}tɐ-jips{]}      & 影子    \\
-ms & {[}kɐ-nə.pjɐ̂ms{]} & 烤火    \\
-ŋs & {[}tɐ-wsoŋs{]}     & 驱邪的香薰\\
   \bottomrule
\end{tabular}
\end{table}

目前观察到，米亚罗话形态过程能够生成的复辅音韵尾有/nɲ/、/mɲ/、/ŋɲ/、/nʃ/、/mʃ/、/ŋʃ/、/kʃ/、/pʃ/、/wj/、/kj/、/js/、/mjs/、/jps/。